\section{What happens when a measure fails}
The theory's distinctive claim is that a coalition requirement prices delay, not denial. This section follows every refused measure forward. Table~\ref{tab:response} collects the accounting; the text walks through it.

% auto-generated by make exhibits -- do not edit
\begin{table}[!htbp]\centering\footnotesize
\caption{The response margin: what refusal buys}\label{tab:response}
\begin{tabular}{@{}l@{\quad}l@{}}
\toprule
Quantity & Value \\
\midrule
\multicolumn{2}{@{}l}{\emph{Panel A: the fate of 100 marginal refusals (2005--19 cohort, $|m|\le 5$, n=422)}} \\ \addlinespace
\quad Re-approved by voters $\le$4y & 54.3 \\
\quad Returned, not yet converted & 13.3 \\
\quad Issued via board or statutory channel & 5.2 \\
\quad Issued on pre-existing voter authority & 9.0 \\
\quad Extinguished within horizon & 18.2 \\
\addlinespace \multicolumn{2}{@{}l}{\emph{First post-vote authoriser ($|m|\le5$): barely-passed vs barely-failed}} \\
\quad barely-PASSED: voter 48.4\%, board 7.9\%, none 43.1\% & \\
\quad barely-FAILED: voter 36.7\%, board 8.6\%, none 54.2\% & \\
\addlinespace \multicolumn{2}{@{}l}{\emph{Panel B: re-submission (2,680 failed GO measures)}} \\ \addlinespace
\quad Hazard, years 1/2/3/4 (\%) & 26.7 / 22.8 / 15.8 / 12.4 \\
\quad Cumulative return within 4y (\%) & 58.2 \\
\quad Median time to return (years) & 1.02 \\
\quad Returns that pass (\%) & 61.9 \\
\quad Median amount ratio, return/original & 1.000 \\
\quad Purpose categories retained on return (\%) & 78.2 \\
\bottomrule
\end{tabular}
\begin{minipage}{0.96\linewidth}\vspace{2pt}\scriptsize
\emph{Notes:} Fates are mutually exclusive per 100 barely-refused measures; issuance on pre-existing voter authority draws on authorisations banked before the refused measure and is not evidence of evading it. Panel B: full failed sample; amounts where registries print them (n=1,354). Retention from the audited purpose bridge (precision 80.0\%).
\end{minipage}
% BEGIN READING (strip for submission)
\begin{minipage}{0.96\linewidth}\vspace{4pt}\scriptsize
\textbf{Reading.} Refusal is a pause: most marginal refusals are re-approved at the next election at the same ask, fewer than one in five extinguish, and the board channel is a floor rather than the treatment margin.
\end{minipage}
% END READING
\end{table}

\textbf{Refusal is a pause, not a verdict.} Of 2,680 failed GO measures, well over half return to the ballot within four years, most at the very next election, and when they return they usually win. Districts do not come back chastened: the median returning measure asks for exactly the original amount, and four-fifths of the original purposes survive into the re-submission. Adding up the fates of 100 barely-refused measures: 54 are re-approved by voters within four years, 13 more have returned and await a verdict, 9 issue on authorisations their voters had approved earlier (banked authority, not evasion), only 5 borrow through a board or statutory channel, and 18 are extinguished within the horizon. The transition matrix beneath Table~\ref{tab:response} makes the same point from the other side: what the cutoff moves is voter-authorised issuance against no issuance at all, while board-channel issuance sits near 8\% on both sides of the line. For the menu-poor sector, the board channel is a floor, not an escape route.

\textbf{The project itself usually survives.} Matching each measure's stated purpose to the classified purposes of the government's subsequent bonds, a third of narrowly refused projects are financed within six years anyway, against 45\% of narrowly passed ones, arriving over a year later. Refusal costs time and forces a second campaign; it only rarely kills the school extension. And the pay-go bound of Section 6 closes the remaining loophole: the gap is not being filled by debt-free construction.

\textbf{One exception: supermajority near-misses bank consent that stalls.} Within California, measures that won a majority but missed the supermajority threshold behave differently. They re-submit far more often than decisive failures and convert to passage at four times the rate, yet the authorisations they eventually win sit undrawn: median time from re-vote to first issue is over six years, against under three for ordinary passes, and their issuance deficit widens rather than closes between years six and eight (Table~\ref{tab:chain}). Where the bar is a supermajority, even re-assembled consent does not become borrowing on the study horizon. A candidate mechanism, statutory tax-rate caps metering the drawdown, could be neither supported nor excluded on the available cells; the fact stands documented and unexplained.

